%=====================================================================
% Datos · Relaciones entre tablas
%=====================================================================
\subsection*{Relaciones entre tablas}
\addcontentsline{toc}{subsection}{Relaciones entre tablas}

Cada consulta de esta subsección une tablas por sus llaves foráneas. Si un registro apuntara a otro que no existe, la unión lo perdería y la consulta devolvería menos filas de las que tiene la tabla; la base, además, no lo permite, como muestra la prueba 5 de los casos que la base rechaza.

\subsubsection*{La partida y sus participantes}
\addcontentsline{toc}{subsubsection}{La partida y sus participantes}

\textit{Partida} se relaciona con \textit{Modo} (N:1) y con \textit{Usuario} a través de \textit{Participacion}, la entidad asociativa que guarda el resultado y el elo de cada jugador. La consulta devuelve las \dpNumPartidas{} partidas con sus participantes: dos en las de dos jugadores, cuatro en la de cuatro y uno en la de la IA, que no es un usuario.

\consulta{relpartidas}
% Archivo generado por bd/exportar_datos.py. No editar a mano.
\begin{center}\scriptsize\setlength{\tabcolsep}{2pt}
\begin{longtable}{|L{35.0pt}|L{82.4pt}|L{77.6pt}|L{60.8pt}|L{44.2pt}|L{64.0pt}|L{44.7pt}|L{27.1pt}|}
\hline
\textbf{\hspace{0pt}Partida} & \textbf{\hspace{0pt}Tipo} & \textbf{\hspace{0pt}Modo} & \textbf{\hspace{0pt}Estado} & \textbf{\hspace{0pt}Jugador} & \textbf{\hspace{0pt}Cuenta} & \textbf{\hspace{0pt}Resultado} & \textbf{\hspace{0pt}Elo final} \\ \hline
\endfirsthead
\hline
\textbf{\hspace{0pt}Partida} & \textbf{\hspace{0pt}Tipo} & \textbf{\hspace{0pt}Modo} & \textbf{\hspace{0pt}Estado} & \textbf{\hspace{0pt}Jugador} & \textbf{\hspace{0pt}Cuenta} & \textbf{\hspace{0pt}Resultado} & \textbf{\hspace{0pt}Elo final} \\ \hline
\endhead
1 & CLASIFICATORIA & CLASICO & FINALIZADA & luis\_\allowbreak{}quo & REGISTRADA & GANADA & 1016 \\ \hline
1 & CLASIFICATORIA & CLASICO & FINALIZADA & marta\_\allowbreak{}muros & REGISTRADA & PERDIDA & 984 \\ \hline
2 & CLASIFICATORIA & CLASICO & FINALIZADA & marta\_\allowbreak{}muros & REGISTRADA & PERDIDA & 969 \\ \hline
2 & CLASIFICATORIA & CLASICO & FINALIZADA & luis\_\allowbreak{}quo & REGISTRADA & GANADA & 1031 \\ \hline
3 & CLASIFICATORIA & CLASICO & FINALIZADA & luis\_\allowbreak{}quo & REGISTRADA & PERDIDA & 1012 \\ \hline
3 & CLASIFICATORIA & CLASICO & FINALIZADA & marta\_\allowbreak{}muros & REGISTRADA & GANADA & 988 \\ \hline
4 & CLASIFICATORIA & CLASICO & FINALIZADA & marta\_\allowbreak{}muros & REGISTRADA & PERDIDA & 973 \\ \hline
4 & CLASIFICATORIA & CLASICO & FINALIZADA & luis\_\allowbreak{}quo & REGISTRADA & GANADA & 1027 \\ \hline
5 & CLASIFICATORIA & CLASICO & FINALIZADA & luis\_\allowbreak{}quo & REGISTRADA & GANADA & 1041 \\ \hline
5 & CLASIFICATORIA & CLASICO & FINALIZADA & marta\_\allowbreak{}muros & REGISTRADA & PERDIDA & 959 \\ \hline
6 & CLASIFICATORIA & RAPIDA & FINALIZADA & pedro\_\allowbreak{}peon & REGISTRADA & TABLAS & 1000 \\ \hline
6 & CLASIFICATORIA & RAPIDA & FINALIZADA & sofia\_\allowbreak{}salto & REGISTRADA & TABLAS & 1000 \\ \hline
7 & CLASIFICATORIA & CUATRO\_\allowbreak{}JUGADORES & FINALIZADA & luis\_\allowbreak{}quo & REGISTRADA & GANADA & 1024 \\ \hline
7 & CLASIFICATORIA & CUATRO\_\allowbreak{}JUGADORES & FINALIZADA & marta\_\allowbreak{}muros & REGISTRADA & PERDIDA & 1008 \\ \hline
7 & CLASIFICATORIA & CUATRO\_\allowbreak{}JUGADORES & FINALIZADA & pedro\_\allowbreak{}peon & REGISTRADA & PERDIDA & 976 \\ \hline
7 & CLASIFICATORIA & CUATRO\_\allowbreak{}JUGADORES & FINALIZADA & sofia\_\allowbreak{}salto & REGISTRADA & PERDIDA & 992 \\ \hline
8 & IA & CLASICO & FINALIZADA & pedro\_\allowbreak{}peon & REGISTRADA & PERDIDA & \textit{NULL} \\ \hline
9 & CLASIFICATORIA & CLASICO & FINALIZADA & luis\_\allowbreak{}quo & REGISTRADA & GANADA & 1055 \\ \hline
9 & CLASIFICATORIA & CLASICO & FINALIZADA & rayo\_\allowbreak{}veloz & REGISTRADA & PERDIDA & 986 \\ \hline
10 & PRIVADA & RAPIDA & FINALIZADA & sofia\_\allowbreak{}salto & REGISTRADA & GANADA & \textit{NULL} \\ \hline
10 & PRIVADA & RAPIDA & FINALIZADA & Invitado\_\allowbreak{}4821 & INVITADO & PERDIDA & \textit{NULL} \\ \hline
11 & CLASIFICATORIA & RAPIDA & FINALIZADA & marta\_\allowbreak{}muros & REGISTRADA & GANADA & 1016 \\ \hline
11 & CLASIFICATORIA & RAPIDA & FINALIZADA & pedro\_\allowbreak{}peon & REGISTRADA & PERDIDA & 984 \\ \hline
12 & CLASIFICATORIA & CLASICO & EN\_\allowbreak{}CURSO & luis\_\allowbreak{}quo & REGISTRADA & \textit{NULL} & \textit{NULL} \\ \hline
12 & CLASIFICATORIA & CLASICO & EN\_\allowbreak{}CURSO & sofia\_\allowbreak{}salto & REGISTRADA & \textit{NULL} & \textit{NULL} \\ \hline
\end{longtable}
\end{center}


\captura{04-relaciones}{Captura 3. La misma consulta en la consola: veinticinco participaciones de doce partidas.}

\subsubsection*{Del jugador a su aspecto}
\addcontentsline{toc}{subsubsection}{Del jugador a su aspecto}

\textit{Equipamiento} tiene dos llaves foráneas compuestas: (\at{id\_usuario}, \at{id\_objeto}) hacia \textit{UsuarioObjeto}, que obliga a equipar solo lo que se posee (\mbox{CU-14} \mbox{RN-02}), e (\at{id\_objeto}, \at{id\_ranura}) hacia \textit{ObjetoCosmetico}, que obliga a que el objeto sea de esa ranura (\mbox{D-03}). La consulta recorre la cadena completa para \at{luis\_quo}, con una fila por ranura:

\consulta{relequipamiento}
% Archivo generado por bd/exportar_datos.py. No editar a mano.
\begin{center}\footnotesize\setlength{\tabcolsep}{3pt}
\begin{longtable}{|L{0.118\linewidth}|L{0.103\linewidth}|L{0.265\linewidth}|L{0.074\linewidth}|L{0.235\linewidth}|L{0.103\linewidth}|}
\hline
\textbf{Usuario} & \textbf{Ranura} & \textbf{Objeto} & \textbf{Rareza} & \textbf{Obtenido} & \textbf{Origen} \\ \hline
\endfirsthead
\hline
\textbf{Usuario} & \textbf{Ranura} & \textbf{Objeto} & \textbf{Rareza} & \textbf{Obtenido} & \textbf{Origen} \\ \hline
\endhead
luis\_\allowbreak{}quo & PEON & PEON\_\allowbreak{}CRISTAL & RARO & 2026-07-27 18:00 & COMPRA \\ \hline
luis\_\allowbreak{}quo & MURO & MURO\_\allowbreak{}MADERA & COMUN & 2026-06-10 16:00 & INICIAL \\ \hline
luis\_\allowbreak{}quo & TABLERO & TABLERO\_\allowbreak{}CLASICO & COMUN & 2026-06-10 16:00 & INICIAL \\ \hline
luis\_\allowbreak{}quo & TITULO & TITULO\_\allowbreak{}CONSTRUCTOR & RARO & 2026-07-20 19:10 & NIVEL \\ \hline
luis\_\allowbreak{}quo & MARCO & MARCO\_\allowbreak{}SENCILLO & COMUN & 2026-06-10 16:00 & INICIAL \\ \hline
luis\_\allowbreak{}quo & EMOTES & EMOTE\_\allowbreak{}RISA & RARO & 2026-07-21 18:00 & COMPRA \\ \hline
\end{longtable}
\end{center}


Los seis objetos llegan por tres caminos distintos: los iniciales del alta (\mbox{CU-02} \mbox{RN-08}), el título que ganó al subir de nivel y los dos que compró.

\subsubsection*{Del reporte a la apelación}
\addcontentsline{toc}{subsubsection}{Del reporte a la apelación}

Un \textit{Reporte} tiene dos relaciones con \textit{Usuario}, una por papel, y una con \textit{MotivoReporte}; una \textit{Sancion} puede venir de un reporte, y una \textit{Apelacion} pertenece a una sanción. Las uniones externas conservan los reportes sin sanción y las sanciones sin apelación:

\consulta{relmoderacion}
% Archivo generado por bd/exportar_datos.py. No editar a mano.
\begin{center}\scriptsize\setlength{\tabcolsep}{3pt}
\begin{longtable}{|L{0.045\linewidth}|L{0.151\linewidth}|L{0.087\linewidth}|L{0.103\linewidth}|L{0.159\linewidth}|L{0.045\linewidth}|L{0.135\linewidth}|L{0.057\linewidth}|L{0.072\linewidth}|}
\hline
\textbf{Reporte} & \textbf{Motivo} & \textbf{Denunciante} & \textbf{Reportado} & \textbf{Estado del reporte} & \textbf{Sanción} & \textbf{Tipo} & \textbf{Apelación} & \textbf{Estado de la apelación} \\ \hline
\endfirsthead
\hline
\textbf{Reporte} & \textbf{Motivo} & \textbf{Denunciante} & \textbf{Reportado} & \textbf{Estado del reporte} & \textbf{Sanción} & \textbf{Tipo} & \textbf{Apelación} & \textbf{Estado de la apelación} \\ \hline
\endhead
1 & LENGUAJE\_\allowbreak{}OFENSIVO & luis\_\allowbreak{}quo & rayo\_\allowbreak{}veloz & RESUELTO\_\allowbreak{}CON\_\allowbreak{}SANCION & 1 & CUENTA TEMPORAL & \textit{NULL} & \textit{NULL} \\ \hline
1 & LENGUAJE\_\allowbreak{}OFENSIVO & luis\_\allowbreak{}quo & rayo\_\allowbreak{}veloz & RESUELTO\_\allowbreak{}CON\_\allowbreak{}SANCION & 2 & CUENTA TEMPORAL & 2 & PENDIENTE \\ \hline
2 & ACOSO & marta\_\allowbreak{}muros & rayo\_\allowbreak{}veloz & PENDIENTE & \textit{NULL} & \textit{NULL} & \textit{NULL} & \textit{NULL} \\ \hline
3 & JUEGO\_\allowbreak{}ANTIDEPORTIVO & sofia\_\allowbreak{}salto & luis\_\allowbreak{}quo & EN\_\allowbreak{}REVISION & \textit{NULL} & \textit{NULL} & \textit{NULL} & \textit{NULL} \\ \hline
4 & NOMBRE\_\allowbreak{}INAPROPIADO & luis\_\allowbreak{}quo & 4dm1n\_\allowbreak{}oficial & RESUELTO\_\allowbreak{}CON\_\allowbreak{}SANCION & 5 & CUENTA PERMANENTE & \textit{NULL} & \textit{NULL} \\ \hline
5 & TRAMPAS & marta\_\allowbreak{}muros & luis\_\allowbreak{}quo & RESUELTO\_\allowbreak{}SIN\_\allowbreak{}SANCION & \textit{NULL} & \textit{NULL} & \textit{NULL} & \textit{NULL} \\ \hline
\end{longtable}
\end{center}


El reporte 1 produjo dos sanciones: la segunda sustituye a la primera, y es la que rayo apeló. El 4 terminó en el baneo de \at{4dm1n\_oficial}; el 5 se resolvió sin sanción, y el 2 y el 3 siguen abiertos.

\subsubsection*{De la caja a su contenido}
\addcontentsline{toc}{subsubsection}{De la caja a su contenido}

\textit{Caja} pertenece a una cuenta y a un \textit{TipoCaja}; \textit{CajaContenido} es su entidad débil, con los tres objetos que salieron, y cada objeto repetido genera un \textit{MovimientoMoneda} de conversión que apunta a la caja y al objeto:

\consulta{relcajas}
% Archivo generado por bd/exportar_datos.py. No editar a mano.
\begin{center}\scriptsize\setlength{\tabcolsep}{3pt}
\begin{longtable}{|L{0.041\linewidth}|L{0.140\linewidth}|L{0.140\linewidth}|L{0.115\linewidth}|L{0.061\linewidth}|L{0.166\linewidth}|L{0.102\linewidth}|L{0.102\linewidth}|}
\hline
\textbf{Caja} & \textbf{Usuario} & \textbf{Tipo} & \textbf{Estado} & \textbf{Número} & \textbf{Objeto} & \textbf{Monedas por conversión} & \textbf{Movimiento} \\ \hline
\endfirsthead
\hline
\textbf{Caja} & \textbf{Usuario} & \textbf{Tipo} & \textbf{Estado} & \textbf{Número} & \textbf{Objeto} & \textbf{Monedas por conversión} & \textbf{Movimiento} \\ \hline
\endhead
1 & marta\_\allowbreak{}muros & CAJA\_\allowbreak{}MADERA & CADUCADA & \textit{NULL} & \textit{NULL} & \textit{NULL} & \textit{NULL} \\ \hline
2 & pedro\_\allowbreak{}peon & CAJA\_\allowbreak{}MADERA & SIN\_\allowbreak{}ABRIR & \textit{NULL} & \textit{NULL} & \textit{NULL} & \textit{NULL} \\ \hline
3 & luis\_\allowbreak{}quo & CAJA\_\allowbreak{}MADERA & ABIERTA & 1 & MURO\_\allowbreak{}LADRILLO & \textit{NULL} & \textit{NULL} \\ \hline
3 & luis\_\allowbreak{}quo & CAJA\_\allowbreak{}MADERA & ABIERTA & 2 & EMOTE\_\allowbreak{}RISA & 40 & 28 \\ \hline
3 & luis\_\allowbreak{}quo & CAJA\_\allowbreak{}MADERA & ABIERTA & 3 & MURO\_\allowbreak{}HIELO & \textit{NULL} & \textit{NULL} \\ \hline
\end{longtable}
\end{center}


Solo la caja abierta tiene contenido. El \at{EMOTE\_RISA} que salió en ella ya lo tenía luis, así que se convirtió en 40 monedas (\mbox{CU-39} \mbox{RN-03}) con el movimiento 28.

\subsubsection*{Una sanción que sustituye a otra}
\addcontentsline{toc}{subsubsection}{Una sanción que sustituye a otra}

La única relación de una tabla consigo misma: \at{id\_sancion\_sustituta} de \textit{Sancion} apunta a otra fila de \textit{Sancion} (\mbox{CU-44} \mbox{RN-09}). La sanción sustituida queda retirada, con su fecha de retiro, y no se borra:

\consulta{relsustitucion}
% Archivo generado por bd/exportar_datos.py. No editar a mano.
\begin{center}\footnotesize\setlength{\tabcolsep}{3pt}
\begin{longtable}{|L{0.055\linewidth}|L{0.139\linewidth}|L{0.139\linewidth}|L{0.062\linewidth}|L{0.139\linewidth}|L{0.364\linewidth}|}
\hline
\textbf{Sanción retirada} & \textbf{Fin previsto} & \textbf{Retiro} & \textbf{Sustituta} & \textbf{Fin de la sustituta} & \textbf{Motivo de la sustituta} \\ \hline
\endfirsthead
\hline
\textbf{Sanción retirada} & \textbf{Fin previsto} & \textbf{Retiro} & \textbf{Sustituta} & \textbf{Fin de la sustituta} & \textbf{Motivo de la sustituta} \\ \hline
\endhead
1 & 2026-08-28 10:00 & 2026-08-27 16:00 & 2 & 2026-09-03 16:00 & Insultos reiterados; se amplía la suspensión. \\ \hline
\end{longtable}
\end{center}

